% =========================================================================
% Chapter 3 - Fourier Neural Operators
% =========================================================================
\chapter{Fourier neural operators}
\label{ch:fno}

\section{Learning operators, not functions}

A conventional network learns a map between finite-dimensional vectors. Retrain
it for a different input size and nothing transfers. That is fatal here: the
grid $(N_1,N_2,N_3)$ of a plane-wave calculation is fixed by the cell and the
cutoff, so \emph{every material has a different one}.

A \emph{neural operator} instead learns a map between function spaces,
\begin{equation}
  \mathcal{G}: \mathcal{A} \longrightarrow \mathcal{U},
  \qquad
  a(\rr) \longmapsto u(\rr),
\end{equation}
and is then discretised only for evaluation. Its parameters describe the
operator, not a particular sampling of it, so one set of weights serves every
grid.

\section{Kernel integral operators}

The construction generalises the affine layer. Lift the input pointwise to a
higher-dimensional channel space, $v_0(\rr) = P\,a(\rr)$, apply $L$ layers of
the form
\begin{equation}
  v_{\ell+1}(\rr) = \sigma\!\left(
      W v_{\ell}(\rr) + b
      + \bigl(\mathcal{K}(v_\ell)\bigr)(\rr) \right),
  \label{eq:noperator}
\end{equation}
and project back pointwise, $u(\rr) = Q\,v_L(\rr)$. Here $W$ is a local
(channel-mixing) linear map and $\mathcal{K}$ is a \emph{kernel integral
operator}
\begin{equation}
  \bigl(\mathcal{K}v\bigr)(\rr) = \int_D \kappa(\rr,\rr')\,v(\rr')\,\mathrm{d}\rr' .
\end{equation}
The pointwise terms alone would give a network with no spatial coupling at all;
$\mathcal{K}$ is what makes it an operator.

\section{The Fourier layer}

Restricting the kernel to be translation invariant, $\kappa(\rr,\rr')
= \kappa(\rr-\rr')$, turns the integral into a convolution, which the
convolution theorem diagonalises:
\begin{equation}
  \bigl(\mathcal{K}v\bigr)(\rr)
    = \Fop^{-1}\!\left[ R(\GG)\cdot \Fop[v](\GG) \right](\rr),
  \label{eq:fourierlayer}
\end{equation}
where $R(\GG)$ is a learned complex multiplier. Rather than parameterising
$\kappa$ in real space, the Fourier neural operator parameterises $R$ directly
and \emph{truncates} it to the lowest $M_1{\times}M_2{\times}M_3$ modes.

Truncation does three things at once: it regularises (high frequencies are
where noise lives), it bounds the parameter count, and --- decisively --- it
makes the weight tensor independent of the grid. Nothing in $R$ refers to
$(N_1,N_2,N_3)$.

\subsection{Why this suits a crystal}

Three properties align the architecture with the physics of
Chapters~\ref{ch:ksdft}--\ref{ch:ofdft}.

\begin{description}
  \item[Periodicity is exact, not learned.] The FFT imposes Born--von Kármán
    boundary conditions, which a crystal unit cell already obeys. A padded
    convolutional network would have to spend capacity imitating that, and
    would still produce edge artefacts.
  \item[The coupling is global in one layer.] Multiplying in reciprocal space
    is convolving over the whole cell in real space. This directly addresses
    the non-locality of Section~\ref{sec:hk}: the Hartree kernel
    $4\pi e^2/G^2$ that the operator must partly represent is
    \emph{diagonal} in exactly the basis the layer works in.
  \item[Differential operators are free and exact.] On a plane-wave grid
    $\nabla \to i\GG$ and $\nabla^2 \to -G^2$ are exact for band-limited
    fields. Since the layer already computes FFTs, the physics-informed terms
    of Chapter~\ref{ch:pifno} cost almost nothing --- and carry no
    discretisation error that the loss would otherwise misattribute to the
    network.
\end{description}

\section{Discretisation invariance}
\label{sec:discinv}

For a real input the transform is taken with \texttt{rfftn}, halving the last
axis. The retained coefficients form four ``corners'' in the first two axes
--- combinations of low positive and low negative frequencies --- each with its
own weight block. The parameter count is
\begin{equation}
  4 \times C_{\mathrm{in}} \times C_{\mathrm{out}} \times M_1 M_2 M_3
  \quad\text{complex numbers,}
\end{equation}
independent of the grid.

Two further choices are required for weights to be \emph{meaningful} across
resolutions, not merely applicable.

\subsection{Normalisation}

The discrete transform must approximate the continuous Fourier series. With
coefficients
\begin{equation}
  c_{\GG} = \frac{1}{N}\sum_{n} u_n e^{-i\GG\cdot\rr_n},
\end{equation}
i.e. \texttt{norm="forward"} in both directions, $c_\GG$ converges to the
continuous series coefficient as $N\to\infty$ and its magnitude does not drift
with $N$. Under the default convention a $120^3$ field would enter a layer with
amplitudes $\sim125\times$ those of a $24^3$ field and shared weights would be
meaningless.

\begin{ptip}
Measured: a band-limited field sampled at $16^3$ and $32^3$ gives operator
outputs agreeing to a relative $1.5\times10^{-6}$ at the shared points.
\end{ptip}

\subsection{Mode selection}

Truncating at a fixed mode \emph{index} keeps a different physical wavevector
band for every cell size, since $|\GG_{\max}| = 2\pi M/L$. Truncating instead
at a fixed $G_{\max}$ --- retaining
\begin{equation}
  M_i = \left\lfloor \frac{G_{\max} L_i}{2\pi} \right\rfloor
\end{equation}
modes --- makes the operator act on the same band of physics in every material.
Poraqu\^e offers both; the physical selection is preferred when the data set
spans a wide range of cell sizes.

\section{Handling grids that differ between materials}

Three mechanisms together deliver true grid independence.

\begin{enumerate}
  \item \textbf{Dynamic mode truncation.} Weights are allocated for $M_i$
    modes; each forward pass uses
    $\min(M_i, \text{modes available on this grid})$ and leaves the rest
    untouched. A $24^3$ and a $120^3$ sample share the same parameters.
  \item \textbf{Resolution-invariant normalisation}, as above.
  \item \textbf{Shape-bucketed batching.} Samples of different spatial shape
    cannot be stacked. Rather than pad --- which wastes compute and injects
    fictitious vacuum into the FFT --- materials are grouped by
    $(N_1,N_2,N_3)$ and batches drawn within a group.
\end{enumerate}

Normalisation layers must also be shape-agnostic: \texttt{GroupNorm} is used
throughout, since batch statistics are meaningless when every sample has a
different spatial extent.

\section{Cell conditioning}

Two materials can share a grid shape and differ entirely in physical scale, so
the lattice must enter explicitly. Poraqu\^e supplies it twice: as three
fractional-coordinate channels appended to the input (bounded in $[0,1)$ for
every material), and through FiLM conditioning, in which a small network maps
rotation-invariant cell descriptors --- the three lattice lengths, the three
angle cosines, and $\Omega^{1/3}$ --- to per-channel scale and shift
parameters. Neither touches the spatial dimensions, so both are oblivious to
the grid.

\section{Architecture summary}

\begin{center}
\begin{tabular}{@{}ll@{}}
\toprule
Stage & Operation \\
\midrule
Input & $(B, 1, N_1, N_2, N_3)$, normalised field \\
Coordinates & append 3 fractional-coordinate channels \\
Lifting & $1\times1\times1$ convolution to width $C$ \\
Fourier layers ($\times L$) & Eq.~\eqref{eq:fourierlayer} + pointwise + GroupNorm + FiLM, residual \\
Projection & two $1\times1\times1$ convolutions with GELU \\
Output & $(B, 1, N_1, N_2, N_3)$ \\
\bottomrule
\end{tabular}
\end{center}

The residual connection around each Fourier layer stabilises deeper stacks; the
pointwise branch supplies the local, high-frequency content that mode
truncation removes from the spectral branch.
